% concatenated from Automorphic_de_Branges_-_Rovnyak_07-RD.tex
%%%%%%%%%%%%%%%%%%%%%%%%%%%%%%%%%%%%%%%%%%%%%%%%%%%%%%%%%%
%\documentclass[11pt, draft]{amsart}
\documentclass[reqno,11pt,centertags]{article}
\usepackage{amsmath,amsthm,amscd,amssymb,latexsym,upref}
\date{\today}
%%%%%%%%%%%%%%%%%%%%%%%%%%%%%%%%%%%%%%%%

%%%\input epsf
%%%\usepackage{epsfig}
%\usepackage{psfig}
\usepackage[T2A,OT1]{fontenc}
%\usepackage[ot2enc]{inputenc}
\usepackage[russian,english]{babel}
\usepackage[margin=3.5 cm]{geometry}
%%%\usepackage{epic,eepicemu}

\usepackage{hyperref}
\hypersetup{
    colorlinks=true, %set true if you want colored links
    linktoc=all,     %set to all if you want both sections and subsections linked
    linkcolor=blue,  %choose some color if you want links to stand out
}

%\usepackage{relsize}

%\setlength{\textwidth}{125mm}
%\setlength{\textheight}{185mm}
%\setlength{\parindent}{8mm}
%\setlength{\oddsidemargin}{0pt}
%\setlength{\evensidemargin}{0pt}
%\frenchspacing
%%%%%%%%%%%%%%%%%%%%%%%%%%%%%%%%%%%%%%%%%%%%%%%%%%%%%%%%%%%
%HERE you TURN ON/OFF the tags for eqs., refs., etc.%
%%%%  \usepackage{showkeys}
%%%%%%%%%%%%%%%%%%%%%%%%%%%%%%%%%%%%%%%%%%%%%%%%%%%%%%%%%%%
\newcommand{\bbD}{{\mathbb{D}}}
\newcommand{\bbE}{{\mathbb{E}}}
\newcommand{\bbN}{{\mathbb{N}}}
\newcommand{\bbR}{{\mathbb{R}}}
\newcommand{\bbP}{{\mathbb{P}}}
\newcommand{\bbZ}{{\mathbb{Z}}}
\newcommand{\bbC}{{\mathbb{C}}}
\newcommand{\bbQ}{{\mathbb{Q}}}
\newcommand{\bbT}{{\mathbb{T}}}
\newcommand{\bbI}{{\mathbb{I}}}

\newcommand{\cA}{{\mathcal{A}}}
\newcommand{\cB}{{\mathcal{B}}}
\newcommand{\cC}{{\mathcal{C}}}
\newcommand{\cD}{{\mathcal{D}}}
\newcommand{\cE}{{\mathcal{E}}}
\newcommand{\cF}{{\mathcal{F}}}
\newcommand{\cG}{{\mathcal{G}}}
\newcommand{\cH}{{\mathcal{H}}}
\newcommand{\cJ}{{\mathcal{J}}}
\newcommand{\cL}{{\mathcal{L}}}
\newcommand{\cM}{{\mathcal{M}}}
\newcommand{\cO}{{\mathcal{O}}}
\newcommand{\cP}{{\mathcal{P}}}
\newcommand{\cQ}{{\mathcal{Q}}}
\newcommand{\cR}{{\mathcal{R}}}
\newcommand{\cS}{{\mathcal{S}}}
\newcommand{\cZ}{{\mathcal{Z}}}
\newcommand{\cw}{{\mathfcal{w}}}

\newcommand{\fA}{{\mathfrak{A}}}
\newcommand{\fB}{{\mathfrak{B}}}
\newcommand{\fH}{{\mathfrak{H}}}
\newcommand{\fF}{{\mathfrak{F}}}
\newcommand{\fK}{{\mathfrak{K}}}
\newcommand{\fa}{{\mathfrak{a}}}
\newcommand{\fb}{{\mathfrak{b}}}
\newcommand{\fe}{{\mathfrak{e}}}
\newcommand{\ff}{{\mathfrak{f}}}
\newcommand{\fj}{{\mathfrak{j}}}
\newcommand{\fu}{{\mathfrak{u}}}
\newcommand{\fv}{{\mathfrak{v}}}
\newcommand{\fw}{{\mathfrak{w}}}
\newcommand{\fz}{{\mathfrak{z}}}

\newcommand{\bC}{{\mathbf{C}}}
\newcommand{\ba}{{\mathbf{a}}}
\newcommand{\bb}{{\mathbf{b}}}
\newcommand{\bc}{{\mathbf{c}}}
\newcommand{\bk}{{\mathbf{k}}}

\newcommand{\e}{{\epsilon}}
\renewcommand{\k}{\varkappa}
\newcommand{\om}{{\omega}}
\newcommand{\z}{\zeta}

\newcommand{\vp}{{\vec{p}}}
\newcommand{\vq}{{\vec{q}}}

%%%%%%%%%%%%%%%%%%%%%%%%%%%%%%%%%%%%%%%%%%%%%%%%%%%%%%%%%
\newcommand{\oL}{\circ\!\!\!L}
\newcommand{\oc}{\overset{\circ}}
\newcommand{\is}{\mathcal{IS}_E}
\renewcommand{\Re}{\text{\rm Re}\,}
\renewcommand{\Im}{\text{\rm Im}\,}
\newcommand{\GSMP}{\text{\rm GSMP}}
\newcommand{\KS}{\text{\rm KS}}
\newcommand{\KSA}{\text{\rm KSA}}
\newcommand{\tr}{\text{\rm tr}\,}
\newcommand{\Diag}{\text{\rm Diag}}
\newcommand{\dist}{\text{\rm dist}}
\newcommand{\Res}{\text{\rm Res}\,}
\newcommand{\sgn}{\text{\rm sgn}}
\newcommand{\Mth}{\langle \theta \rangle_{I}}
%%%%%%%%%%%%%%%%%%%%%%%%%%%%%%%%%%%%%%%%%%%
\allowdisplaybreaks \numberwithin{equation}{section}
%%%%%%%%%%%%%%%%%%%%%%%%%%%%%%%%%%%%%%%%%%%%%%
\newtheorem{theorem}{Theorem}[section]


\newtheorem{theoremm}{Theorem}[section]
\newtheorem{lemma}[theorem]{Lemma}
\newtheorem{proposition}[theorem]{Proposition}
\newtheorem{claim}[theorem]{Claim}
\newtheorem{corollary}[theorem]{Corollary}
\newtheorem{hypothesis}[theorem]{Hypothesis}
\theoremstyle{definition}
\newtheorem{definition}[theorem]{Definition}
\newtheorem{notation}[theorem]{Notation}
\newtheorem{question}[theorem]{Question}
\newtheorem{remark}[theorem]{Remark}
\newtheorem{problem}[theorem]{Problem}
\newtheorem{fact}[theorem]{Fact}
\newtheorem{example}[theorem]{Example}
\newtheorem{summary}[theorem]{Summary}


%%%%%%%%%%%%%%%%% Add. Macros %%%%%%%%%
\def\be{\begin{equation}}
\def\ee{\end{equation}}
\def\bea{\begin{eqnarray}}
\def\eea{\end{eqnarray}}
\def\bean{\begin{eqnarray*}}
\def\eean{\end{eqnarray*}}
\def\bbm{\begin{bmatrix}}
\def\ebm{\end{bmatrix}}
\def\lla{\left\langle}
\def\rra{\right\rangle}


\def\sep{\;\vrule\;}
\def\restr#1{\,\vrule\,\lower1ex\hbox{$#1$}}

\def\a{\alpha}
\def\b{\beta}
\def\d{\delta}
\def\D{\Delta}
\def\e{\epsilon}
\def\f{\varphi}
\def\F{\Phi}
\def\g{\gamma}
\def\G{\Gamma}
\def\i{\infty}
\def\k{\kappa}
\def\l{\lambda}
\def\L{\Lambda}
\def\o{\omega}
\def\O{\Omega}
\def\r{\rho}
\def\s{\sigma}
\def\S{\Sigma}
\def\t{\theta}
\def\z{\zeta}
\def\R{{\bf R}}
\def\N{{\bf N}}
\def\C{{\bf C}}
\def\c{{\rm cap}}
\def\li#1{\lim_{#1\to\i}}

\def\setm{\setminus}
%%%%%%%%%%%%%%%%%%%%%%%%%%%%

%\date{\today}
\title
{Automorphic de Branges - Rovnyak Spaces}
\author{Alexander Kheifets}
%

\begin{document}

\maketitle

\begin{center}
\textit{Dedicated to the memory of Professor Heinz Langer}
\end{center}
\medskip


%% \tableofcontents



\begin{abstract}

A family of automorphic de Branges - Rovnyak
spaces (depending on an arbitrary character $\a$)
is associated to a $\b$-automorphic analytic function $w$
bounded in modulus by $1$. Multiplication by the Green function
defines unitary operators acting between those spaces. 
Function $w$ can be recovered from these operators as the characteristic function.
It is shown that for Nevanlinna-Pick problem this family of unitary operators extends
the family of isometric operators defined by the interpolation data.



\end{abstract}

\section{Automorphic Hardy Spaces}

\subsection{Green Function. Widom. Pommerenke.}

Let $\Gamma$ be a  Fuchsian group acting on the unit disk $\mathbb D$.
A mapping $\a$ from a group $\G$ to the unit circle $\bbT$ is called a character of $\G$ if
$$
\a(\g_1\g_2)=\a(\g_1)\a(\g_2).
$$
We say that function $f$ defined on the unit disk $\mathbb D$ or/and on the unit circle $\mathbb T$
is character automorphic (more specifically $\a$ automorphic) if
$$
f\circ\g=\a(\g) f,\quad\forall \g\in\G.
$$
Assuming that $\G$ is of convergent type, we define
the Green function of $\Gamma$ at $0$ as
$$
g_0(\zeta)=\prod\limits_{\gamma\in\Gamma}\gamma(\zeta)\frac{|\gamma(0)|}{\gamma(0)}.
$$
Group $\G$ is said to be of Widom type \cite{Widom71, Pom} if $g'_0$ is of bounded characteristic. In this case
\begin{theorem}[Widom-Pommerenke]\label{June19-02}
$$
g'_0(\zeta)=\dfrac{\Delta_0(\zeta)}{\psi_0(\zeta)},
$$
where $\Delta_0$ is an inner function 
\footnote{We will always use normalization $\D_0(0)>0$.}
and $\psi_0$ is a an outer function bounded in modulus
by~$1$.
\end{theorem}
Note that $g_0$ and $\D_0$ are character automorphic functions. We denote their characters as $\nu_0$
and $\d_0$, respectively.
\if{%%%%%%%%%%%%%%%%%%%%%%%%%%%%%%%%%%%%%%%%%%%%%%%%%%%%%%%%%%%%%%%%%%%%%%%%%%%%%%%%%%%%%%%%%%
We will also consider $g_{\zeta_0}(\zeta)$ the Green function of $\Gamma$ at $\zeta_0$
$$
g_{\zeta_0}(\zeta)=
\prod\limits_{\gamma\in\Gamma}\frac{\gamma(\zeta)-\zeta_0}{1-\gamma(\zeta)\overline{\zeta_0}}
c_{\gamma, \zeta_0}.
$$
For groups of Widom type we also have
$$
g'_{\z_0}(\zeta)=\dfrac{\Delta_{\z_0}(\zeta)}{\psi_{\z_0}(\zeta)},
$$
where $\Delta_{\z_0}$ is an inner function and $\psi_{\z_0}$ is a bounded outer function. We will use notations
$\nu_{\z_0}$ for the character of $g_{\z_0}$ and $\d_{\z_0}$ for the character of $\D_{\z_0}$.
}\fi%%%%%%%%%%%%%%%%%%%%%%%%%%%%%%%%%%%%%%%%%%%%%%%%%%%%%%%%%%%%%%%%%%%%%%%%%%%%%%%%%%%%%%%%%%%

\subsection{Automorphic Hardy Spaces. Reproducing Kernels. Direct Cauchy Theorem.}\label{260111-01}

For groups of Widom type all spaces $H^p(\a)$, $1\le p\le\infty$, $\a\in\G^*$, contain nonconstant functions.
This can be seen by means of the Poincar\' e theta series
$$
(P^\alpha f) (t)=\dfrac{\sum\limits_{\gamma\in\Gamma}\overline{\alpha(\gamma)}f(\gamma(t))
|\gamma'(t)|}
{\sum\limits_{\gamma\in\Gamma}|\gamma'(t)|}
$$
$$
=\frac{\sum\limits_{\gamma\in\Gamma}\overline{\alpha(\gamma)}f(\gamma(t))
\dfrac{\gamma'(t)}{\g(t)}}
{\sum\limits_{\gamma\in\Gamma}\dfrac{\gamma'(t)}{\g(t)}}
=\frac{\sum\limits_{\gamma\in\Gamma}\overline{\alpha(\gamma)}f(\gamma(t))
\dfrac{\gamma'(t)}{\g(t)}}
{\dfrac{g'_0(t)}{g_0(t)}}.
$$
$P^\a$ is a contractive projection from $L^p$ onto $L^p(\a)$. For $p=2$, $P^\a$
is also selfadjoint (that is, it is the orthogonal projection from $L^2$ onto $L^2(\a)$). When applied to functions of
the form $\D_0 H^p$ it produces many functions in $H^p(\a)$. However, $P^\a$ is not the orthogonal projection from $H^2$
onto $H^2(\a)$.

By $k^\a_\z(t)$ we denote the orthogonal projection of $\dfrac{1}{1-t\overline\z}$ on $H^2(\a)$. It is the reproducing kernel of $H^2(\a)$ at point $\z$.

We say that the Direct Cauchy Theorem holds for a group of Widom type $\G$
(see \cite{Hasu, SY})
if for every $f\in H^1(\d_0)$ the following equality holds
\be\label{240901-01}
\int\limits_{\mathbb T}\frac{f}{\D_0}L(dt)=\frac{f(0)}{\D_0(0)},
\ee
where $L(dt)$ is the normalized Lebesgue measure on the unit circle $\bbT$. In particular, \eqref{240901-01} implies
that
$$
k^{\d_0}_0=\D_0\overline{\D_0(0)}
$$
is the reproducing kernel at $0$ for $H^2(\d_0)$.

The following orthogonal decomposition of
$L^2(\a)$ is an equivalent form of the Direct Cauchy Theorem property
(see \cite{Hasu, SY})
\be\label{250302-02}
H^2_\perp(\alpha)=L^2(\alpha)\ominus H^2(\alpha)
=\frac{\Delta_0}{g_0}\overline{H^2(\delta_0\overline{\nu_0}\ \overline\alpha)}.
\ee
Let $\widetilde k^\alpha_\zeta\in H^2_\perp(\alpha)$ be such that
$$
\left\langle
\dfrac{\D_0}{g_0}\overline{h}, \widetilde k^\alpha_\zeta
\right\rangle_{H^2_\perp(\alpha)}=\overline{h(\zeta)}.
$$
It is straightforward that
\be\label{June22-01}
\widetilde k^\alpha_\zeta
=\dfrac{\Delta_0}{g_0}\overline{k^{\delta_0 \overline{\nu_0}\ \overline\alpha}_\zeta}
.
\ee
The next lemma is a combination of Lemma 7.2 in \cite{SY}
and a weaker version of Lemma~3.2 in \cite{VY}.
We prove it here for the reader's convenience.
\begin{lemma}\label{250607_01}
Assume that the Direct Cauchy Theorem holds for group $\G$. Then
\be\label{250608_03}
\D_0\dfrac{\overline{k^\alpha_0}}{\sqrt{k^\alpha_0(0)}}
 =
\dfrac{k^{\d_0\overline\a}_0}{\sqrt{k^{\d_0\overline\a}_0(0)}}
\ee
and
\be\label{250608_04}
\D_0(0)
=\sqrt{k^\alpha_0(0)k^{\d_0\overline\a}_0 (0)}.
\ee
\end{lemma}
\begin{proof}
First we show that
$\D_0\overline{k^{\a}_0}\perp H^2_\perp(\d_0\overline\alpha)$,
that is,
$\D_0\overline{k^{\a}_0}\in H^2(\d_0\overline\alpha)$. Indeed,
by formula \eqref{250302-02},
$$
H^2_\perp(\d_0\overline\alpha)
=\frac{\Delta_0}{g_0}\overline{H^2(\overline{\nu_0}\ \alpha)}.
$$
Then for every $h\in {H^2(\overline{\nu_0}\ \alpha)}$
$$
\left\langle
\dfrac{\D_0}{g_0}\overline{h}, \D_0\overline{k^\alpha_0}
\right\rangle
=
\left\langle
k^\alpha_0 , g_0 h
\right\rangle
=0.
$$
Next, for every $u\in {H^2(\d_0\ \overline\alpha)}$
$$
\left\langle
u, \D_0\overline{k^\alpha_0}
\right\rangle
=
\int\limits_{\bbT}\dfrac{u(t)k^\alpha_0(t)}{\D_0(t)}L(dt)
=
\dfrac{u(0)k^\alpha_0(0)}{\D_0(0)},
$$
by the Direct Cauchy Theorem property. Therefore (since $\D_0(0)>0$),
\be\label{250608_01}
\D_0\overline{k^\alpha_0}
 =
\dfrac{k^\alpha_0(0)}{\D_0(0)}k^{\d_0\overline\a}_0.
\ee
Taking squares of the norms on the both sides of \eqref{250608_01} we get
$$
k^\alpha_0(0)
 =
\dfrac{k^\alpha_0(0)^2}{\D_0(0)^2}k^{\d_0\overline\a}_0 (0),
$$
that is,
\be\label{250608_02}
\D_0(0)
=
\sqrt{k^\alpha_0(0)k^{\d_0\overline\a}_0 (0)}.
\ee
Substituting \eqref{250608_02} into \eqref{250608_01} we get \eqref{250608_03}.
\end{proof}



\section{Automorphic de Branges-Rovnyak spaces}

%%\subsection{Automorphic de Branges-Rovnyak Spaces of the Disk}

%%%The classical de Branges - Rovmyak space is defined as follows (see, e.g., \cite{Kh}).
Let $L^2$ be the space of square summable functions against the Lebesgue measure
 on the unit circle $\bbT$.
Let $w$ be an analytic, bounded in modulus by $1$ function on $\bbD$ and let
$W(t):=
\begin{bmatrix}
1& w(t) \\
\overline{w(t)} & 1 
\end{bmatrix}
$. 
The space $L^w$ is the range
space $W^{1/2}(L_2\oplus L_2)$ endowed with the range norm. In more
detail: for every element $f$ in $L^w$, there exists a unique $h_f\in
L_2\oplus L_2$ which is orthogonal to ${\rm Ker}\, W(t)$ for almost all
$t\in\mathbb T$ and such that $f=W^{1/2}h_f$. This unique $h_f$ will be denoted
by $h_f:=W^{[-1/2]}f$ and the $L^w$-norm of $f$ is defined as
$$
\|f\|_{L^w}^2:=\|h_{_f}\|_{L_2\oplus L_2}^2=
\int_{\mathbb T}\left\|
%%{\scriptsize\left[\begin{array}{cc}
\begin{bmatrix}
1& w(t) \\
\overline{w(t)} & 1
\end{bmatrix}
%%\end{array}\right]}
^{[-1/2]}f(t)\right\|^2_{\C^2}L(dt),
$$
where $L(dt)$ stands for the normalized Lebesgue
measure on $\mathbb T$. The classical de Branges~-~Rovnyak space $H^w$
is defined (see \cite{BR}) as a subspace of $L^w$ that consists of functions
$$
f=\begin{bmatrix} f_2 \\ f_1\end{bmatrix},\quad f_2\in H^2_+,\quad f_1\in H^2_-,
$$
where $H^2_+$ and $H^2_-$ are the standard Hardy subspaces of $L^2$.


Assume that $w$ is character automorphic with character $\beta$. Note that then
$$
W\circ\gamma=
\begin{bmatrix}
1 & w\circ\gamma \\
\overline{w\circ\gamma} & 1
\end{bmatrix}
%%%%%%%%=
%%%%%%%%\begin{bmatrix}
%%%%%%%%\sqrt{\beta(\gamma)} & 0 \\
%%%%%%%%0 & \overline{\sqrt{\beta(\gamma)}}
%%%%%%%%\end{bmatrix}
%%%%%%%%\begin{bmatrix}
%%%%%%%%1 & w \\
%%%%%%%%\overline{w} & 1
%%%%%%%%\end{bmatrix}
%%%%%%%%\begin{bmatrix}
%%%%%%%%\overline{\sqrt{\beta(\gamma)}} & 0
%%%%%%%%\\
%%%%%%%%0 & \sqrt{\beta(\gamma)}
%%%%%%%%\end{bmatrix}
%%%%%%%%$$
%%%%%%%%$$
=
\begin{bmatrix}
\beta(\gamma) & 0 \\
0 & 1
\end{bmatrix}
\begin{bmatrix}
1 & w \\
\overline{w} & 1
\end{bmatrix}
\begin{bmatrix}
\overline{\beta(\gamma)} & 0
\\
0 & 1
\end{bmatrix}
=
\begin{bmatrix}
\beta(\gamma) & 0 \\
0 & 1
\end{bmatrix}
W
\begin{bmatrix}
\overline{\beta(\gamma)} & 0
\\
0 & 1
\end{bmatrix}
$$
for every $\g\in\G$.
The same relation holds for $W^{1/2}$. We consider now subspace $L^w(\alpha, \beta)$ of $L^w$
\be\label{250302-01}
\begin{bmatrix} f_2 \\ f_1 \end{bmatrix}\in L^w,\quad f_2\in L^2(\alpha\beta), \quad f_1\in L^2(\alpha),
\ee
where $L^2(\alpha)$ is a subspace of $L^2$ that consists of $\a$ automorphic functions.
To motivate: functions of the form
$$
\begin{bmatrix} f_2 \\ f_1 \end{bmatrix}=
W^{1/2}\begin{bmatrix} h_2 \\ h_1 \end{bmatrix},\quad h_2\in L^2(\alpha\beta), \quad h_1\in L^2(\alpha)
$$
satisfy \eqref{250302-01}. We define $H^w(\alpha, \beta)$ as a subspace of $L^w(\alpha, \beta)$ such that
$$
f_2\in H^2_+(\alpha\beta),\quad f_1\in H^2_\perp(\alpha)=
\dfrac{\D_0}{g_0}\overline{H^2(\overline{\nu_0}\d_0\overline\alpha)}
.
$$
Thus, to every $\b$-automorphic analytic function $w$, bounded in modulus by $1$,
one can associate a family of subspaces $L^w(\a,\b)$ and $H^w(\a,\b)$ that depend on
arbitrary character~$\a$.



\section{Associated Unitary Operators and Unitary Colligations}
Using \eqref{250302-02} one can get the following orthogonal decomposition
ot space $L^w(\a,\b)$
$$
L^w(\a,\b)=
\begin{bmatrix}
1 & w \\
\overline{w} & 1
\end{bmatrix}
\begin{bmatrix}
0
\\
H^2(\a)
\end{bmatrix}
\oplus
H^w(\a, \b)
\oplus
\begin{bmatrix}
1 & w \\
\overline{w} & 1
\end{bmatrix}
\begin{bmatrix}
H^2_{\perp}(\a\b)
\\
0
\end{bmatrix}
$$
$$
=
\begin{bmatrix}
1 & w \\
\overline{w} & 1
\end{bmatrix}
\begin{bmatrix}
0
\\
H^2(\a)
\end{bmatrix}
\oplus
H^w(\a, \b)
\oplus
\begin{bmatrix}
1 & w \\
\overline{w} & 1
\end{bmatrix}
\begin{bmatrix}
\dfrac{\D_0}{g_0}\overline{H^2(\overline{\nu_0}\d_0\overline{\alpha\b})}
\\
0
\end{bmatrix}.
$$
The operator of multiplication by $\overline{g_0}$ maps unitarily $L^w(\a,\b)$
onto $L^w(\overline{\nu_0}\a,\b)$.
Thus we have a family of unitary operators acting on the family of spaces $L^w(\a,\b)$.
Since any unitary operator preserves orthgonality, the
multiplication by $\overline{g_0}$  also maps unitarily
$$
\left\{
\dfrac{1}{\sqrt{k^\a_0(0)}}
\begin{bmatrix}
1 & w \\
\overline{w} & 1
\end{bmatrix}
\begin{bmatrix}
0
\\
k^\a_0
\end{bmatrix}
\right\}
\oplus
H^w(\a, \b)
$$
onto
$$
H^w(\overline{\nu_0}\a, \b)
\oplus
\left\{
\dfrac{1}{\sqrt{k^{\d_0\overline{\alpha\b}}_0(0)}}
\begin{bmatrix}
1 & w \\
\overline{w} & 1
\end{bmatrix}
\begin{bmatrix}
\dfrac{\D_0}{g_0}\overline{k^{\d_0\overline{\alpha\b}}_0}
\\
0
\end{bmatrix}
\right\}
$$
\be\label{241225-01}%%%%%%%%%%%%%%%%%%%{240907-01}
=
H^w(\overline{\nu_0}\a, \b)
\oplus
\left\{
\dfrac{1}{\sqrt{k^{\alpha\b}_0(0)}}
\begin{bmatrix}
1 & w \\
\overline{w} & 1
\end{bmatrix}
\begin{bmatrix}
\overline{g_0}k^{\alpha\b}_0
\\
0
\end{bmatrix}
\right\}
.
\ee
The latter equality in \eqref{241225-01} %%%%%%%%%%%%%%%%%%{240907-01} 
is due to Lemma \ref{250607_01}.
Thus we get a family of unitary colligations. 
%%%%%%%Sometimes it is convenient to
%%%%%%%have normalized vectors in the braces
%%%%%%%$$
%%%%%%%\overline{g_0}:\
%%%%%%%\left\{
%%%%%%%\dfrac{1}{\sqrt{k^\a_0(0)}}
%%%%%%%\begin{bmatrix}
%%%%%%%1 & w \\
%%%%%%%\overline{w} & 1
%%%%%%%\end{bmatrix}
%%%%%%%\begin{bmatrix}
%%%%%%%0
%%%%%%%\\
%%%%%%%k^\a_0
%%%%%%%\end{bmatrix}
%%%%%%%\right\}
%%%%%%%\oplus
%%%%%%%H^w(\a, \b)
%%%%%%%$$
%%%%%%%unitarily onto
%%%%%%%\be\label{241225-01}
%%%%%%%H^w(\overline{\nu_0}\a, \b)
%%%%%%%\oplus
%%%%%%%\left\{
%%%%%%%\dfrac{1}{\sqrt{k^{\alpha\b}_0(0)}}
%%%%%%%\begin{bmatrix}
%%%%%%%1 & w \\
%%%%%%%\overline{w} & 1
%%%%%%%\end{bmatrix}
%%%%%%%\begin{bmatrix}
%%%%%%%\overline{g_0}k^{\alpha\b}_0
%%%%%%%\\
%%%%%%%0
%%%%%%%\end{bmatrix}
%%%%%%%\right\}
%%%%%%%.
%%%%%%%\ee
\begin{remark}
In the next section we show how to compute $w$ from the family of unitary colligations
\eqref{241225-01}, like the characteristic function. However, the author does not know what abstract object is modeled by
this family of colligations.
\end{remark}

\section{$w$ as the Characteristic Function.}

Since $\overline{g_0}$ maps 
$$
H^w(\overline{\nu_0^{n}}\a, \b)
\oplus
\left\{
\dfrac{1}{\sqrt{k^{\overline{\nu_0^{n}}\a}_0(0)}}
\begin{bmatrix}
1 & w \\
\overline{w} & 1
\end{bmatrix}
\begin{bmatrix}
0
\\
k^{\overline{\nu_0^{n}}\a}_0
\end{bmatrix}
\right\}
$$
unitarily onto
%%%%%$$
%%%%%H^w(\overline{\nu^{n+1}_0}\a, \b)
%%%%%\oplus
%%%%%\left\{
%%%%%\dfrac{1}{\sqrt{k^{\d_0\nu^n_0\overline{\alpha\b}}_0(0)}}
%%%%%\begin{bmatrix}
%%%%%1 & w \\
%%%%%\overline{w} & 1
%%%%%\end{bmatrix}
%%%%%\begin{bmatrix}
%%%%%\dfrac{\D_0}{g_0}\overline{k^{\d_0\nu^n_0\overline{\alpha\b}}_0}
%%%%%\\
%%%%%0
%%%%%\end{bmatrix}
%%%%%\right\}
%%%%%$$
$$
%%=
H^w(\overline{\nu^{n+1}_0}\a, \b)
\oplus
\left\{
\dfrac{1}{\sqrt{k^{\overline{\nu^n_0}{\alpha\b}}_0(0)}}
\begin{bmatrix}
1 & w \\
\overline{w} & 1
\end{bmatrix}
\begin{bmatrix}
\overline{g_0} k^{\overline{\nu^n_0}{\alpha\b}}_0
\\
0
\end{bmatrix}
\right\}
,
$$
we consider this dynamics
\be\label{260225_03}
\overline{g_0}
\left(
h_n
\oplus
e_{1.n}
\right)
=
h_{n+1}
\oplus
e_{2,n},\quad n\ge 0,
\ee
where $h_n\in H^w(\overline{\nu_0^{n}}\a, \b)$,
\be\label{260225_04}
e_{1,n}:=
\dfrac{c_{1,n}}{\sqrt{k^{\overline{\nu_0^{n}}\a}_0(0)}}
\begin{bmatrix}
1 & w \\
\overline{w} & 1
\end{bmatrix}
\begin{bmatrix}
0
\\
k^{\overline{\nu_0^{n}}\a}_0
\end{bmatrix}
\ee
and
\be\label{260225_05}
e_{2,n}:=
\dfrac{c_{2,n}}{\sqrt{k^{\overline{\nu^n_0}{\alpha\b}}_0(0)}}
\begin{bmatrix}
1 & w \\
\overline{w} & 1
\end{bmatrix}
\begin{bmatrix}
\overline{g_0} k^{\overline{\nu^n_0}{\alpha\b}}_0
\\
0
\end{bmatrix}.
\ee
Note that $h_n$, $e_{1,n}$ and $e_{2,n}$ belong to different spaces for different $n$.
\begin{theorem}
Let $\a$ be an arbitrary character of $\G$.
Let $h_0=0$ and $c_{1,n}$ be an arbitrary $\ell^2$ input sequence.
Let $c_{2,n}$ be the corresponding output sequence from dynamics \eqref{260225_03}-\eqref{260225_05}.
Let
\be\label{260225_01}
\widetilde c_1=
\sum\limits_{n=0}^\infty
c_{1,n}
g^n_0\dfrac{k^{\overline{\nu_0^{n}}\a}_0}{\sqrt{k^{\overline{\nu_0^{n}}\a}_0(0)}}
\ee
and
\be\label{260225_02}
\widetilde c_2=
\sum\limits_{n=0}^\infty
c_{2,n}
g^n_0\dfrac{k^{\overline{\nu_0^{n}}\a\b}_0}{\sqrt{k^{\overline{\nu_0^{n}}\a\b}_0(0)}}
\ee
be the associated Fourier series $($functions in $H^2(\a)$ and in $H^2(\a\b)$, respectively$)$. Then
\be\label{260224_01}
\widetilde c_2=w \widetilde c_1.
\ee
\end{theorem}
\begin{proof}
We start with a special case when $c_{1,0}=1$ and $c_{1,n}=0$ for
$n\ge 1$. $h_0=0$, by assumption. Then for $n=0$ we have
$$
\overline{g_0}
\left(
0
\oplus
\dfrac{1}{\sqrt{k^{\a}_0(0)}}
\begin{bmatrix}
1 & w \\
\overline{w} & 1
\end{bmatrix}
\begin{bmatrix}
0
\\
k^{\a}_0
\end{bmatrix}
\right)
$$
$$
=
h_{1}
\oplus
\dfrac{c_{2,0}}{\sqrt{k^{{\alpha\b}}_0(0)}}
\begin{bmatrix}
1 & w \\
\overline{w} & 1
\end{bmatrix}
\begin{bmatrix}
\overline{g_0} k^{{\alpha\b}}_0
\\
0
\end{bmatrix}
,
$$
where
$$
c_{2,0}=
\left\langle
\dfrac{\overline{g_0}}{\sqrt{k^\a_0(0)}}
\begin{bmatrix}
1 & w \\
\overline{w} & 1
\end{bmatrix}
\begin{bmatrix}
0
\\
k^\a_0
\end{bmatrix}
,
\dfrac{1}{\sqrt{k^{{\alpha\b}}_0(0)}}
\begin{bmatrix}
1 & w \\
\overline{w} & 1
\end{bmatrix}
\begin{bmatrix}
\overline{g_0} k^{{\alpha\b}}_0
\\
0
\end{bmatrix}
\right\rangle
_{L^w(\overline{\nu_0}\a,\b)}
$$
$$
=
\left\langle
\dfrac{wk^\a_0}{\sqrt{k^\a_0(0)}}
,
\dfrac{k^{\alpha\b}_0}{\sqrt{k^{\alpha\b}_0(0)}}
\right\rangle
_{H^2(\a\b)}
=A^\a_0,
$$
where we will use this notations for the expansion of $\dfrac{wk^\a_0}{\sqrt{k^\a_0(0)}}$ with respect to the
standard orthonormal basis in $H^2(\a\b)$ (see \cite{SY})
\be\label{260117_01}
\dfrac{wk^\a_0}{\sqrt{k^\a_0(0)}}
=
\sum\limits_{n=0}^\infty
A^\a_n
g^n_0\dfrac{k^{\overline{\nu_0^{n}}\a\b}_0}{\sqrt{k^{\overline{\nu_0^{n}}\a\b}_0(0)}}.
\ee
%%%%%%%%$$
%%%%%%%%=
%%%%%%%%\dfrac{1}{\sqrt{k^\a_0(0)}}
%%%%%%%%\dfrac{1}{\sqrt{k^{\d_0\overline{\alpha\b}}_0(0)}}
%%%%%%%%\int\limits_{\bbT}
%%%%%%%%\dfrac
%%%%%%%%{wk^\a_0 k^{\d_0\overline{\alpha\b}}_0}
%%%%%%%%{\D_0}
%%%%%%%%L(dt)
%%%%%%%%$$
%%%%%%%%using the Direct Cauchy Theorem
%%%%%%%%$$
%%%%%%%%=
%%%%%%%%\dfrac{1}{\sqrt{k^\a_0(0)}}
%%%%%%%%\dfrac{1}{\sqrt{k^{\d_0\overline{\alpha\b}}_0(0)}}
%%%%%%%%\dfrac
%%%%%%%%{w(0)k^\a_0(0) k^{\d_0\overline{\alpha\b}}_0(0)}
%%%%%%%%{\D_0(0)}
%%%%%%%%$$
%%%%%%%%$$
%%%%%%%%=
%%%%%%%%w(0)
%%%%%%%%\dfrac
%%%%%%%%{\sqrt{k^\a_0(0)} \sqrt{k^{\d_0\overline{\alpha\b}}_0(0)}}
%%%%%%%%{\D_0(0)}
%%%%%%%%=
%%%%%%%%w(0)
%%%%%%%%\dfrac
%%%%%%%%{\sqrt{k^\a_0(0)} }
%%%%%%%%{\sqrt{k^{{\alpha\b}}_0(0)}}.
%%%%%%%%$$
%%%%%%%%The latter equality is due to Lemma \ref{250607_01}. 
Thus, we get
$$
e_{2,0}=
A^\a_0
\overline{g_0}
\begin{bmatrix}
1 & w \\
\overline{w} & 1
\end{bmatrix}
\begin{bmatrix}
\dfrac{k^{{\alpha\b}}_0}{\sqrt{k^{{\alpha\b}}_0(0)}}
\\ \\
0
\end{bmatrix}
$$
and
$$
h_1=
\overline{g_0}
\begin{bmatrix}
1 & w \\
\overline{w} & 1
\end{bmatrix}
\begin{bmatrix}
-A^\a_0\dfrac{k^{{\alpha\b}}_0}{\sqrt{k^{{\alpha\b}}_0(0)}}
\\ \\
\dfrac{k^\a_0}{\sqrt{k^\a_0(0)}}
\end{bmatrix}.
%%%%%%%%%-
%%%%%%%%%w(0)
%%%%%%%%%\dfrac
%%%%%%%%%{\sqrt{k^\a_0(0)} }
%%%%%%%%%{\sqrt{k^{{\alpha\b}}_0(0)}}
%%%%%%%%%\dfrac{1}{\sqrt{k^{\d_0\overline{\alpha\b}}_0(0)}}
%%%%%%%%%\begin{bmatrix}
%%%%%%%%%1 & w \\
%%%%%%%%%\overline{w} & 1
%%%%%%%%%\end{bmatrix}
%%%%%%%%%\begin{bmatrix}
%%%%%%%%%\dfrac{\D_0}{g_0}\overline{k^{\d_0\overline{\alpha\b}}_0}
%%%%%%%%%\\
%%%%%%%%%0
%%%%%%%%%\end{bmatrix}
$$
Next, for $n=1$, we have
$$
\overline{g_0}
\left(
h_1
\oplus
0
\right)
$$
$$
=
h_{2}
\oplus
\dfrac{c_{2,1}}{\sqrt{k^{\overline{\nu_0}\alpha\b}_0(0)}}
\begin{bmatrix}
1 & w \\
\overline{w} & 1
\end{bmatrix}
\begin{bmatrix}
\overline{g_0} k^{\overline{\nu_0}\alpha\b}_0
\\
0
\end{bmatrix}
,
$$
where
$$
c_{2,1}=
\left\langle
\overline{g_0}h_1
,
\dfrac{1}{\sqrt{k^{\overline{\nu_0}\alpha\b}_0(0)}}
\begin{bmatrix}
1 & w \\
\overline{w} & 1
\end{bmatrix}
\begin{bmatrix}
\overline{g_0} k^{\overline{\nu_0}\alpha\b}_0
\\
0
\end{bmatrix}
\right\rangle
$$
$$
=
\left\langle
\overline{g_0}^2
\begin{bmatrix}
1 & w \\
\overline{w} & 1
\end{bmatrix}
\begin{bmatrix}
-A^\a_0\dfrac{k^{{\alpha\b}}_0}{\sqrt{k^{{\alpha\b}}_0(0)}}
\\ \\
\dfrac{k^\a_0}{\sqrt{k^\a_0(0)}}
\end{bmatrix}
,
\overline{g_0}
\begin{bmatrix}
1 & w \\
\overline{w} & 1
\end{bmatrix}
\begin{bmatrix}
\dfrac{k^{\overline{\nu_0}\alpha\b}_0}{\sqrt{k^{\overline{\nu_0}\alpha\b}_0(0)}}
\\ \\
0
\end{bmatrix}
\right\rangle
_{L^w(\overline{\nu_0}\a,\b)}
$$
$$
=
\left\langle
\left(
\dfrac{wk^\a_0}{\sqrt{k^\a_0(0)}}
-A^\a_0\dfrac{k^{{\alpha\b}}_0}{\sqrt{k^{{\alpha\b}}_0(0)}}
\right)
,
g_0
\dfrac{k^{\overline{\nu_0}\alpha\b}_0}{\sqrt{k^{\overline{\nu_0}\alpha\b}_0(0)}}
\right\rangle
_{H^2(\a\b)}
=A^\a_1.
$$
Thus, we get
$$
e_{2,1}=
A^\a_1
\overline{g_0}
\begin{bmatrix}
1 & w \\
\overline{w} & 1
\end{bmatrix}
\begin{bmatrix}
\dfrac{k^{\overline{\nu_0}\alpha\b}_0}{\sqrt{k^{\overline{\nu_0}\alpha\b}_0(0)}}
\\ \\
0
\end{bmatrix}
$$
and
$$
h_2=
\overline{g_0}^2
\begin{bmatrix}
1 & w \\
\overline{w} & 1
\end{bmatrix}
\begin{bmatrix}
-A^\a_0\dfrac{k^{{\alpha\b}}_0}{\sqrt{k^{{\alpha\b}}_0(0)}}
-A^\a_1 g_0\dfrac{k^{\overline{\nu_0}\alpha\b}_0}{\sqrt{k^{\overline{\nu_0}\alpha\b}_0(0)}}
\\ \\
\dfrac{k^\a_0}{\sqrt{k^\a_0(0)}}
\end{bmatrix}.
$$
Proceeding by induction, we get
$$
c_{2,n}=A^\a_n,
$$
$$
e_{2,n}=
A^\a_n
\overline{g_0}
\begin{bmatrix}
1 & w \\
\overline{w} & 1
\end{bmatrix}
\begin{bmatrix}
\dfrac{k^{\overline{\nu_0}^n\alpha\b}_0}{\sqrt{k^{\overline{\nu_0}^n\alpha\b}_0(0)}}
\\ \\
0
\end{bmatrix}
$$
and (for $n\ge 1$)
$$
h_n=
\overline{g_0}^n
\begin{bmatrix}
1 & w \\
\overline{w} & 1
\end{bmatrix}
\begin{bmatrix}
-A^\a_0\dfrac{k^{{\alpha\b}}_0}{\sqrt{k^{{\alpha\b}}_0(0)}}
-A^\a_1 g_0\dfrac{k^{\overline{\nu_0}\alpha\b}_0}{\sqrt{k^{\overline{\nu_0}\alpha\b}_0(0)}}
-\ldots
-A^\a_{n-1} g_0^{n-1}\dfrac{k^{\overline{\nu_0}^{n-1}\alpha\b}_0}{\sqrt{k^{\overline{\nu_0}^{n-1}\alpha\b}_0(0)}}
\\ \\
\dfrac{k^\a_0}{\sqrt{k^\a_0(0)}}
\end{bmatrix}.
$$
Therefore, the Fourier series \eqref{260225_01} and \eqref{260225_02} will be
$$
\widetilde c_1=
\sum\limits_{n=0}^\infty
c_{1,n}
g^n_0\dfrac{k^{\overline{\nu_0^{n}}\a}_0}{\sqrt{k^{\overline{\nu_0^{n}}\a}_0(0)}}
=
\dfrac{k^\a_0}{\sqrt{k^\a_0(0)}}.
$$
and
$$
\widetilde c_2=
\sum\limits_{n=0}^\infty
c_{2,n}
g^n_0\dfrac{k^{\overline{\nu_0^{n}}\a\b}_0}{\sqrt{k^{\overline{\nu_0^{n}}\a\b}_0(0)}}
=
\sum\limits_{n=0}^\infty
A^\a_n
g^n_0\dfrac{k^{\overline{\nu_0^{n}}\a\b}_0}{\sqrt{k^{\overline{\nu_0^{n}}\a\b}_0(0)}}
=
\dfrac{wk^\a_0}{\sqrt{k^\a_0(0)}}.
$$
Hence \eqref{260224_01} holds for the input sequence $c_{1,0}=1$, $c_{1,n}=0$ for $n\ge 1$.

Let now $c_{1,1}=1$ and $c_{1,n}=0$ for all other indices $n$. Then $h_1=0$, $c_{2,0}=0$ and
similar computation to the above starts with $c_{1,1}=1$ and $h_1=0\in H^w(\overline{\nu_0}\a, \b)$.
Therefore,
$$
c_{2,n}=A^{\overline{\nu_0}\a}_{n-1}, \quad n\ge 1 .
$$
Hence in this case
$$
\widetilde c_1=
\sum\limits_{n=0}^\infty
c_{1,n}
g^n_0\dfrac{k^{\overline{\nu_0^{n}}\a}_0}{\sqrt{k^{\overline{\nu_0^{n}}\a}_0(0)}}
=
g_0\dfrac{k^{\overline{\nu_0}\a}_0}{\sqrt{k^{\overline{\nu_0}\a}_0(0)}}.
$$
and
$$
\widetilde c_2
=
\sum\limits_{n=0}^\infty
c_{2,n}
g^n_0\dfrac{k^{\overline{\nu_0^{n}}\a\b}_0}{\sqrt{k^{\overline{\nu_0^{n}}\a\b}_0(0)}}
$$
$$
=
\sum\limits_{n=1}^\infty
A^{\overline{\nu_0}\a}_{n-1}
g^n_0\dfrac{k^{\overline{\nu_0^{n}}\a\b}_0}{\sqrt{k^{\overline{\nu_0^{n}}\a\b}_0(0)}}
=
\sum\limits_{n=0}^\infty
A^{\overline{\nu_0}\a}_{n}
g^{n+1}_0\dfrac{k^{\overline{\nu_0^{n+1}}\a\b}_0}{\sqrt{k^{\overline{\nu_0^{n+1}}\a\b}_0(0)}}
$$
$$
=
g_0\sum\limits_{n=0}^\infty
A^{\overline{\nu_0}\a}_{n}
g^{n}_0\dfrac{k^{\overline{\nu_0^{n+1}}\a\b}_0}{\sqrt{k^{\overline{\nu_0^{n+1}}\a\b}_0(0)}}
=
g_0
\dfrac{wk^{\overline{\nu_0}\a}_0}{\sqrt{k^{\overline{\nu_0}\a}_0(0)}}.
$$
Therefore, \eqref{260224_01} holds for this $c_1$ as well. Similar computation shows that
\eqref{260224_01} holds for every standard basic vector $c_1$ in $\ell^2$. Finally, due to the linearity of
dynamics \eqref{260225_03}-\eqref{260225_05},
\eqref{260224_01} holds for every $c_1\in\ell^2$.
\end{proof}



%%%%\medskip
%%%%{\bf Questions}
%%%%
%%%%\begin{enumerate}
%%%%\item What does the family of unitary operators \eqref{241225-01} model?
%%%%\item
%%%%In what sense is $w$ the characteristic function of the family?
%%%%\end{enumerate}

%%%%%%%%%%%%%%\section{Notations ?}
%%%%%%%%%%%%%%Let
%%%%%%%%%%%%%%$$
%%%%%%%%%%%%%%L^\a:=
%%%%%%%%%%%%%%L^w(\a, \b),\qquad
%%%%%%%%%%%%%%L^\a_1=
%%%%%%%%%%%%%%\begin{bmatrix}
%%%%%%%%%%%%%%1 & w \\
%%%%%%%%%%%%%%\overline{w} & 1
%%%%%%%%%%%%%%\end{bmatrix}
%%%%%%%%%%%%%%\begin{bmatrix}
%%%%%%%%%%%%%%0
%%%%%%%%%%%%%%\\
%%%%%%%%%%%%%%L^2(\a)
%%%%%%%%%%%%%%\end{bmatrix},\qquad
%%%%%%%%%%%%%%L^\a_2=
%%%%%%%%%%%%%%\begin{bmatrix}
%%%%%%%%%%%%%%1 & w \\
%%%%%%%%%%%%%%\overline{w} & 1
%%%%%%%%%%%%%%\end{bmatrix}
%%%%%%%%%%%%%%\begin{bmatrix}
%%%%%%%%%%%%%%L^2(\a\b)
%%%%%%%%%%%%%%\\
%%%%%%%%%%%%%%0
%%%%%%%%%%%%%%\end{bmatrix},
%%%%%%%%%%%%%%$$
%%%%%%%%%%%%%%$L^\a$ is spanned by $L^\a_2$ and $L^\a_1$.
%%%%%%%%%%%%%%Then
%%%%%%%%%%%%%%$$
%%%%%%%%%%%%%%U^\a=\overline{g_0}:\ L^\a\to L^{\overline{\nu_0}\a},
%%%%%%%%%%%%%%$$
%%%%%%%%%%%%%%and
%%%%%%%%%%%%%%$$
%%%%%%%%%%%%%%U^\a=\overline{g_0}:\ L^\a_i\to L^{\overline{\nu_0}\a}_i
%%%%%%%%%%%%%%$$
%%%%%%%%%%%%%%unitarily onto.
%%%%%%%%%%%%%%Let
%%%%%%%%%%%%%%$$
%%%%%%%%%%%%%%H^\a:=
%%%%%%%%%%%%%%H^w(\a, \b)
%%%%%%%%%%%%%%$$
%%%%%%%%%%%%%%$$
%%%%%%%%%%%%%%E^\a_1:=
%%%%%%%%%%%%%%\left\{
%%%%%%%%%%%%%%\dfrac{1}{\sqrt{k^\a_0(0)}}
%%%%%%%%%%%%%%\begin{bmatrix}
%%%%%%%%%%%%%%1 & w \\
%%%%%%%%%%%%%%\overline{w} & 1
%%%%%%%%%%%%%%\end{bmatrix}
%%%%%%%%%%%%%%\begin{bmatrix}
%%%%%%%%%%%%%%0
%%%%%%%%%%%%%%\\
%%%%%%%%%%%%%%k^\a_0
%%%%%%%%%%%%%%\end{bmatrix}
%%%%%%%%%%%%%%\right\}
%%%%%%%%%%%%%%$$
%%%%%%%%%%%%%%$$
%%%%%%%%%%%%%%E^{\overline{\nu_0}\a}_2 :=
%%%%%%%%%%%%%%\left\{
%%%%%%%%%%%%%%\dfrac{1}{\sqrt{k^{\d_0\overline{\alpha\b}}_0(0)}}
%%%%%%%%%%%%%%\begin{bmatrix}
%%%%%%%%%%%%%%1 & w \\
%%%%%%%%%%%%%%\overline{w} & 1
%%%%%%%%%%%%%%\end{bmatrix}
%%%%%%%%%%%%%%\begin{bmatrix}
%%%%%%%%%%%%%%\dfrac{\D_0}{g_0}\overline{k^{\d_0\overline{\alpha\b}}_0}
%%%%%%%%%%%%%%\\
%%%%%%%%%%%%%%0
%%%%%%%%%%%%%%\end{bmatrix}
%%%%%%%%%%%%%%\right\}
%%%%%%%%%%%%%%.
%%%%%%%%%%%%%%$$
%%%%%%%%%%%%%%Then
%%%%%%%%%%%%%%$$
%%%%%%%%%%%%%%A^\a=\overline{g_0}:\ H^\a\oplus E^\a_1\to
%%%%%%%%%%%%%%H^{\overline{\nu_0}\a}
%%%%%%%%%%%%%%\oplus
%%%%%%%%%%%%%%E^{\overline{\nu_0}\a}_2
%%%%%%%%%%%%%%$$
%%%%%%%%%%%%%%unitarily onto. 
%%%%%%%%%%%%%%Recall that with $n\ge 0$
%%%%%%%%%%%%%%$$
%%%%%%%%%%%%%%g_0^{n}
%%%%%%%%%%%%%%\dfrac
%%%%%%%%%%%%%%{k^{{\nu_0^{-n}}{\a}}_0}
%%%%%%%%%%%%%%{\sqrt{k^{{\nu_0^{-n+1}}{\a}}_0(0)}}
%%%%%%%%%%%%%%$$
%%%%%%%%%%%%%%is an orthonormal basis for $H^2(\a)$ and with $n\le -1$
%%%%%%%%%%%%%%$$
%%%%%%%%%%%%%%g_0^{n}
%%%%%%%%%%%%%%\dfrac
%%%%%%%%%%%%%%{k^{{\nu_0^{-n}}{\a}}_0}
%%%%%%%%%%%%%%{\sqrt{k^{{\nu_0^{-n}}{\a}}_0(0)}}
%%%%%%%%%%%%%%=
%%%%%%%%%%%%%%g_0^{n}\D_0
%%%%%%%%%%%%%%\dfrac
%%%%%%%%%%%%%%{\overline{k^{\d_0\nu_0^{n}\overline{\alpha}}_0}}
%%%%%%%%%%%%%%{\sqrt{k^{\d_0\nu_0^{n}\overline{\alpha}}_0(0)}}
%%%%%%%%%%%%%%$$
%%%%%%%%%%%%%%is an orthonormal basis for $H^2_\perp(\a)$.....
%%%%%%%%%%%%%%
%%%%%%%%%%%%%%
%%%%%%%%%%%%%%
%%%%%%%%%%%%%%
%%%%%%%%%%%%%%
%%%%%%%%%%%%%%
%%%%%%%%%%%%%%Let
%%%%%%%%%%%%%%$$
%%%%%%%%%%%%%%D_{V^\a}:=
%%%%%%%%%%%%%%\left\{
%%%%%%%%%%%%%%\begin{bmatrix}
%%%%%%%%%%%%%%1 & w \\
%%%%%%%%%%%%%%\overline{w} & 1
%%%%%%%%%%%%%%\end{bmatrix}
%%%%%%%%%%%%%%\begin{bmatrix}
%%%%%%%%%%%%%%0
%%%%%%%%%%%%%%\\
%%%%%%%%%%%%%%k^\a_0
%%%%%%%%%%%%%%\end{bmatrix}
%%%%%%%%%%%%%%\dfrac{k^\a_\z(0)}{k^\a_0(0)}
%%%%%%%%%%%%%%\ \overline{w(\zeta)}
%%%%%%%%%%%%%%\oplus
%%%%%%%%%%%%%%K_{\zeta}^{\alpha, \beta}
%%%%%%%%%%%%%%\right\}
%%%%%%%%%%%%%%\subset
%%%%%%%%%%%%%%E^\a_1
%%%%%%%%%%%%%%\oplus
%%%%%%%%%%%%%%\left\{
%%%%%%%%%%%%%%K_{\zeta}^{\alpha, \beta}
%%%%%%%%%%%%%%\right\}.
%%%%%%%%%%%%%%$$
%%%%%%%%%%%%%%$$
%%%%%%%%%%%%%%R_{V^\a}:=
%%%%%%%%%%%%%%\left\{
%%%%%%%%%%%%%%K_{\zeta}^{\overline{\nu_0}\alpha, \beta}
%%%%%%%%%%%%%%\ \overline{g_0(\z)}
%%%%%%%%%%%%%%\oplus
%%%%%%%%%%%%%%\begin{bmatrix}
%%%%%%%%%%%%%%1 & w \\
%%%%%%%%%%%%%%\overline{w} & 1
%%%%%%%%%%%%%%\end{bmatrix}
%%%%%%%%%%%%%%\begin{bmatrix}
%%%%%%%%%%%%%%\dfrac{\D_0}{g_0}\overline{k^{\d_0\overline{\alpha\b}}_0}
%%%%%%%%%%%%%%\\
%%%%%%%%%%%%%%0
%%%%%%%%%%%%%%\end{bmatrix}
%%%%%%%%%%%%%%\dfrac{k^{\a\b}_\z(0)}{\D_0(0)}
%%%%%%%%%%%%%%\right\}
%%%%%%%%%%%%%%\subset
%%%%%%%%%%%%%%E^{\overline{\nu_0}\a}_2
%%%%%%%%%%%%%%\oplus 
%%%%%%%%%%%%%%\left\{K_{\zeta}^{\overline{\nu_0}\alpha, \beta}\right\}
%%%%%%%%%%%%%%.
%%%%%%%%%%%%%%$$
%%%%%%%%%%%%%%$$
%%%%%%%%%%%%%%N^\a_1 :=
%%%%%%%%%%%%%%D^\perp_{V^\a} =
%%%%%%%%%%%%%%\left(
%%%%%%%%%%%%%%E^\a_1
%%%%%%%%%%%%%%\oplus
%%%%%%%%%%%%%%\left\{
%%%%%%%%%%%%%%K_{\zeta}^{\alpha, \beta}
%%%%%%%%%%%%%%\right\}
%%%%%%%%%%%%%%\right)
%%%%%%%%%%%%%%\ominus
%%%%%%%%%%%%%%D_{V^\a}
%%%%%%%%%%%%%%.
%%%%%%%%%%%%%%$$
%%%%%%%%%%%%%%$$
%%%%%%%%%%%%%%N^{\overline{\nu_0}\a}_2 :=
%%%%%%%%%%%%%%R^\perp_{V^\a} =
%%%%%%%%%%%%%%\left(
%%%%%%%%%%%%%%E^{\overline{\nu_0}\a}_2
%%%%%%%%%%%%%%\oplus
%%%%%%%%%%%%%%\left\{
%%%%%%%%%%%%%%K_{\zeta}^{\overline{\nu_0}\alpha, \beta}
%%%%%%%%%%%%%%\right\}
%%%%%%%%%%%%%%\right)
%%%%%%%%%%%%%%\ominus
%%%%%%%%%%%%%%R_{V^\a}
%%%%%%%%%%%%%%.
%%%%%%%%%%%%%%$$
%%%%%%%%%%%%%%$$
%%%%%%%%%%%%%%H^\a_1 :=H^\a \ominus 
%%%%%%%%%%%%%%\left\{
%%%%%%%%%%%%%%K_{\zeta}^{\alpha, \beta}
%%%%%%%%%%%%%%\right\}
%%%%%%%%%%%%%%$$
%%%%%%%%%%%%%%Then
%%%%%%%%%%%%%%$$
%%%%%%%%%%%%%%A^\a :
%%%%%%%%%%%%%%N^\a_1\oplus H^\a_1
%%%%%%%%%%%%%%\to
%%%%%%%%%%%%%%N^{\overline{\nu_0}\a}_2\oplus H^{\overline{\nu_0}\a_1}_1
%%%%%%%%%%%%%%$$



\section{Reproducing Kernels}
In what follows we assume that $w$ is automorphic with character $\beta$.
\begin{lemma}\label{June21-10}
The following functions
\be
K_{\zeta}^{\alpha, \beta}(t)=\left[\begin{array}{c} K_{\zeta,+}^{\alpha, \beta}(t) \\
K_{\zeta,-}^{\alpha, \beta}(t) \end{array}\right]
=\left[\begin{array}{cc} 1& w(t) \\
\overline{w(t)} & 1\end{array}\right]\left[\begin{array}{c} k^{\alpha\beta}_\zeta(t) \\
-\overline{w(\zeta)} k^\alpha_\zeta(t) \end{array}\right]
\label{12.8}
\ee
and
\be
\widetilde{K}_{\zeta}^{\alpha, \beta}(t)=
\left[\begin{array}{c} \widetilde{K}_{\zeta,+}^{\alpha, \beta}(t) \\
\widetilde{K}_{\zeta,-}^{\alpha, \beta}(t) \end{array}\right]
=\left[\begin{array}{cc} 1& w(t) \\
\overline{w(t)} & 1\end{array}\right]\left[\begin{array}{c} -w(\zeta)\widetilde k^{\alpha\beta}_\zeta(t) \\
\widetilde k^\alpha_\zeta(t) \end{array}\right]   \label{12.9}
\ee
are the reproducing kernels in $H^w(\alpha, \beta)$
for the first component at $\zeta$ and for the second component times
$\dfrac{g_0}{\D_0}$ at $\zeta$ $($which is a conjugate analytic function$)$, respectively.
Where $k^\alpha_\zeta$ and $\widetilde k^\alpha_\zeta$ are defined in Section \ref{260111-01}.
\end{lemma}
\begin{proof}
It follows immediately from definition \eqref{12.8} that
$$
K_{\zeta,+}^{\alpha, \beta}(t)=
k^{\alpha\beta}_\zeta(t)-w(t)\overline{w(\zeta)} k^\alpha_\zeta(t)
\in H^2(\alpha\beta).
$$
Clearly
$$
K_{\zeta,-}^{\alpha, \beta}(t)=
\overline{w(t)}k^{\alpha\beta}_\zeta(t)-\overline{w(\zeta)} k^\alpha_\zeta(t)\in L^2(\a).
$$
We show that it is orthogonal to $H^2(\a)$. For $h\in H^2(\a)$
$$
\left\langle h, K_{\zeta,-}^{\alpha, \beta}\right\rangle=
\left\langle h, \overline w k^{\alpha\beta}_\zeta-\overline{w(\zeta)} k^\alpha_\zeta\right\rangle
$$
$$
=
\left\langle w h, k^{\alpha\beta}_\zeta\rangle-w(\z)\langle h, k^\alpha_\zeta\right\rangle
=w(\z) h(\z) - w(\z) h(\z)=0.
$$
Thus, $K_{\zeta,-}^{\alpha, \beta}\in H^2_\perp(\a)$ and, therefore,
$K_{\zeta}^{\alpha, \beta}\in H^w(\alpha, \beta)$.
Now
$$
\left\langle\begin{bmatrix} f_2 \\ f_1 \end{bmatrix}, K_{\zeta}^{\alpha, \beta}\right\rangle_{H^w(\alpha, \beta)}
=
\left\langle\begin{bmatrix} f_2 \\ f_1 \end{bmatrix},
\begin{bmatrix} k^{\alpha\beta}_\zeta \\
-\overline{w(\zeta)} k^\alpha_\zeta \end{bmatrix}
\right\rangle_{L^2}
$$
$$
=
\left\langle f_2 , k^{\alpha\beta}_\zeta \right\rangle_{H^2(\alpha\beta)}
+
\left\langle f_1 , -\overline{w(\zeta)} k^{\alpha}_\zeta \right\rangle_{L^2 (\alpha)}
= f_2(\zeta),
$$
since $f_2\in H^2(\alpha\beta)$ and $f_1\in H^2_\perp(\alpha)$.

Proof of the second part is analogous. In view of \eqref{June22-01},
$$
\widetilde
K_{\zeta,-}^{\alpha, \beta}
=
\widetilde k^\alpha_\zeta(t)
-\overline{w(t)}w(\zeta)\widetilde k^{\alpha\beta}_\zeta(t)
\in
\frac{\D_0}{g_0}\overline{H^2(\overline\nu_0\d_0\overline\alpha)}
=
H^2_\perp(\alpha).
$$
Clearly
$$
\widetilde K_{\zeta,+}^{\alpha, \beta}
=
w(t)\widetilde k^\alpha_\zeta(t)
- w(\zeta)\widetilde k^{\alpha\beta}_\zeta(t)\in L^2(\a\b).
$$
We show that is is orthogonal to $H^2_\perp(\a\b)$.
Since
$$
H^2_\perp(\a\b)=\dfrac{\D_0}{g_0}\overline {H^2(\delta_0 \overline{\nu_0}\ \overline{\alpha\b})}
$$
we compute
$$
\left\langle\widetilde K_{\zeta,+}^{\alpha, \beta}, \dfrac{\D_0}{g_0}\overline h\right\rangle
=\left\langle
w \widetilde k^\alpha_\zeta
- w(\zeta)\widetilde k^{\alpha\beta}_\zeta, \dfrac{\D_0}{g_0}\overline h\right\rangle
$$
(using again \eqref{June22-01})
$$
=\left\langle
w \dfrac{\Delta_0}{g_0}\overline{k^{\delta_0 \overline{\nu_0}\ \overline\alpha}_\zeta}
- w(\zeta)\dfrac{\Delta_0}{g_0}\overline{k^{\delta_0 \overline{\nu_0}\ \overline{\alpha\b}}_\zeta}, \dfrac{\D_0}{g_0}\overline h\right\rangle
$$
$$
=\left\langle
w \overline{k^{\delta_0 \overline{\nu_0}\ \overline\alpha}_\zeta}
- w(\zeta)\overline{k^{\delta_0 \overline{\nu_0}\ \overline{\alpha\b}}_\zeta}, \overline h\right\rangle
$$
$$
=\left\langle
w h, {k^{\delta_0 \overline{\nu_0}\ \overline\alpha}_\zeta}
\right\rangle
-
w(\zeta)
\left\langle h, {k^{\delta_0 \overline{\nu_0}\ \overline{\alpha\b}}_\zeta}\right\rangle
$$
$$
=w(\z)h(\z)-w(\z)h(z)=0.
$$
Now
$$
\left\langle\begin{bmatrix} f_2 \\ f_1 \end{bmatrix}, \widetilde K_{\zeta}^{\alpha, \beta}\right\rangle_{H^w(\alpha, \beta)}
=
\left\langle\begin{bmatrix} f_2 \\ f_1 \end{bmatrix},
\begin{bmatrix} -w(\zeta)\widetilde k^{\alpha\beta}_\zeta \\
\widetilde k^\alpha_\zeta \end{bmatrix}
\right\rangle_{L^2}
$$
$$
=
\left\langle f_2 , -w(\zeta)\widetilde k^{\alpha\beta}_\zeta\right\rangle_{L^2(\alpha\beta)}
+
\left\langle f_1 , \widetilde k^\alpha_\zeta \right\rangle_{H^2_\perp (\alpha)}
=\left(\frac{g_0}{\D_0}f_1\right)(\zeta),
$$
since $f_2\in H^2(\alpha\beta)$ and $f_1\in H^2_\perp(\alpha)$.
\end{proof}


\begin{remark}\label{200524_01}
Since $k_\z^\a\perp H^2_\perp(\a)$, we have that
$$
\left[\begin{array}{cc} 1& w(t) \\
\overline{w(t)} & 1\end{array}\right]\left[\begin{array}{c} 0 \\
k^\alpha_\zeta(t) \end{array}\right]
\perp
H^w(\alpha, \beta).
$$
Then, in view of \eqref{12.8},
$$
K_{\zeta}^{\alpha, \beta}(t)
=\left[\begin{array}{cc} 1& w(t) \\
\overline{w(t)} & 1\end{array}\right]\left[\begin{array}{c} k^{\alpha\beta}_\zeta(t) \\
-\overline{w(\zeta)} k^\alpha_\zeta(t) \end{array}\right]
$$
$$
=P_{H^w(\a,\b)}
\left[\begin{array}{cc} 1& w(t) \\
\overline{w(t)} & 1\end{array}\right]\left[\begin{array}{c} k^{\alpha\beta}_\zeta(t) \\
-\overline{w(\zeta)} k^\alpha_\zeta(t) \end{array}\right]
$$
\be
=P_{H^w(\a,\b)}
\left[\begin{array}{cc} 1& w(t) \\
\overline{w(t)} & 1\end{array}\right]\left[\begin{array}{c} k^{\alpha\beta}_\zeta(t) \\
c k^\alpha_\zeta(t) \end{array}\right],
\label{20524_02}
\ee
where $c$ is any complex number, in particular $0$.
\end{remark}

\section{Towards Interpolation}
If $w$ satisfies certain interpolation conditions, we will have a family of isometric
colligations contained in the family of unitary colligations \eqref{241225-01}.
For example
\begin{lemma}
Multiplication by $\overline{g_0}$ maps
$$
\begin{bmatrix}
1 & w \\
\overline{w} & 1
\end{bmatrix}
\begin{bmatrix}
0
\\
k^\a_0
\end{bmatrix}
\dfrac{k^\a_\z(0)}{k^\a_0(0)}
\ \overline{w(\zeta)}
\oplus
K_{\zeta}^{\alpha, \beta}
$$
to
$$
K_{\zeta}^{\overline{\nu_0}\alpha, \beta}
\ \overline{g_0(\z)}
\oplus
\begin{bmatrix}
1 & w \\
\overline{w} & 1
\end{bmatrix}
\begin{bmatrix}
\dfrac{\D_0}{g_0}\overline{k^{\d_0\overline{\alpha\b}}_0}
\\
0
\end{bmatrix}
\dfrac{k^{\a\b}_\z(0)}{\D_0(0)}
.
$$
\end{lemma}
\begin{proof}
Consider
$$
A
\begin{bmatrix}
1 & w \\
\overline{w} & 1
\end{bmatrix}
\begin{bmatrix}
0
\\
k^\a_0
\end{bmatrix}
\oplus
\left[\begin{array}{cc} 1& w \\
\overline{w} & 1\end{array}\right]\left[\begin{array}{c} k^{\alpha\beta}_\zeta \\
-\overline{w(\zeta)} k^\alpha_\zeta \end{array}\right]
=
\left[\begin{array}{cc} 1& w \\
\overline{w} & 1\end{array}\right]\left[\begin{array}{c} k^{\alpha\beta}_\zeta \\
A k^\a_0-\overline{w(\zeta)} k^\alpha_\zeta \end{array}\right].
$$
We multiply this by $\overline{g_0}$ and compute the projection on
$\begin{bmatrix}
1 & w \\
\overline{w} & 1
\end{bmatrix}
\begin{bmatrix}
\dfrac{\D_0}{g_0}\overline{k^{\d_0\overline{\alpha\b}}_0}
\\
0
\end{bmatrix}.$
To this end we first compute
$$
\left\langle
\overline{g_0}
\left[\begin{array}{cc} 1& w \\
\overline{w} & 1\end{array}\right]\left[\begin{array}{c} k^{\alpha\beta}_\zeta \\
A k^\a_0-\overline{w(\zeta)} k^\alpha_\zeta \end{array}\right]
,
\begin{bmatrix}
1 & w \\
\overline{w} & 1
\end{bmatrix}
\begin{bmatrix}
\dfrac{\D_0}{g_0}\overline{k^{\d_0\overline{\alpha\b}}_0}
\\
0
\end{bmatrix}
\right\rangle_{L^w(\overline{\nu_0}\a.\b)}
$$
$$
=\langle
k^{\alpha\beta}_\zeta +
w(A k^\a_0-\overline{w(\zeta)} k^\alpha_\zeta)
,
\D_0 \overline{k^{\d_0\overline{\alpha\b}}_0}
\rangle_{L^2}
$$
using the Direct Cauchy Theorem
$$
=
\dfrac
{k^{\alpha\beta}_\zeta (0)+
w(0)(A k^\a_0(0)-\overline{w(\zeta)} k^\alpha_\zeta(0))}
{\D_0(0)}
k^{\d_0\overline{\alpha\b}}_0(0)
=
\dfrac
{k^{\alpha\beta}_\zeta (0)}
{\D_0(0)}
k^{\d_0\overline{\alpha\b}}_0(0),
$$
if $A=\overline{w(\zeta)} \dfrac{k^\a_\z(0)}{k^\a_0(0)}$.
Next we need to divide by
$$
\left\langle
\begin{bmatrix}
1 & w \\
\overline{w} & 1
\end{bmatrix}
\begin{bmatrix}
\dfrac{\D_0}{g_0}\overline{k^{\d_0\overline{\alpha\b}}_0}
\\
0
\end{bmatrix}
,
\begin{bmatrix}
1 & w \\
\overline{w} & 1
\end{bmatrix}
\begin{bmatrix}
\dfrac{\D_0}{g_0}\overline{k^{\d_0\overline{\alpha\b}}_0}
\\
0
\end{bmatrix}
\right\rangle_{L^w(\overline{\nu_0}\a.\b)}
=k^{\d_0\overline{\alpha\b}}_0(0).
$$
Thus, the requisite projection is
$$
\begin{bmatrix}
1 & w \\
\overline{w} & 1
\end{bmatrix}
\begin{bmatrix}
\dfrac{\D_0}{g_0}\overline{k^{\d_0\overline{\alpha\b}}_0}
\\
0
\end{bmatrix}
\dfrac{k^{\alpha\beta}_\zeta (0)}{\D_0(0)}
.
$$
Now we consider
$$
\overline{g_0}
\left[\begin{array}{cc} 1& w \\
\overline{w} & 1\end{array}\right]\left[\begin{array}{c} k^{\alpha\beta}_\zeta \\
\overline{w(\zeta)} \dfrac{k^\a_0}{k^\a_0(0)}k^\alpha_\zeta(0)-\overline{w(\zeta)} k^\alpha_\zeta \end{array}\right]
-
\dfrac
{k^{\alpha\beta}_\zeta (0)}
{\D_0(0)}
\begin{bmatrix}
1 & w \\
\overline{w} & 1
\end{bmatrix}
\begin{bmatrix}
\dfrac{\D_0}{g_0}\overline{k^{\d_0\overline{\alpha\b}}_0}
\\
0
\end{bmatrix}
$$
\be\label{240908-02}
=
\left[\begin{array}{cc} 1& w \\
\overline{w} & 1\end{array}\right]\left[\begin{array}{c}
k^{\alpha\beta}_\zeta\overline{g_0}
-
\dfrac{k^{\alpha\beta}_\zeta (0)}{\D_0(0)}
\dfrac{\D_0}{g_0}\overline{k^{\d_0\overline{\alpha\b}}_0}
 \\
\left(
\overline{w(\zeta)} \dfrac{k^\a_0}{k^\a_0(0)}k^\alpha_\zeta(0)-\overline{w(\zeta)} k^\alpha_\zeta
\right)
\overline{g_0}
-\overline w
\dfrac{k^{\alpha\beta}_\zeta (0)}{\D_0(0)}
\dfrac{\D_0}{g_0}\overline{k^{\d_0\overline{\alpha\b}}_0}
\end{array}\right]
\ee
Next we show that
\be\label{240908_01}
k^{\alpha\beta}_\zeta\overline{g_0}
-
\dfrac{k^{\alpha\beta}_\zeta (0)}{\D_0(0)}
\dfrac{\D_0}{g_0}\overline{k^{\d_0\overline{\alpha\b}}_0}
=\overline{g_0(\z)}k^{\overline{\nu_0}\alpha\beta}_\zeta .
\ee
Clearly this function is in $L^2(\overline{\nu_0}\alpha\beta)$. To show that it is in
$H^2(\overline{\nu_0}\alpha\beta)$ we prove that it is orthogonal to
$H^2_\perp(\overline{\nu_0}\alpha\beta)$. Indeed,
$$
\left\langle
k^{\alpha\beta}_\zeta\overline{g_0}
-
\dfrac{k^{\alpha\beta}_\zeta (0)}{\D_0(0)}
\dfrac{\D_0}{g_0}\overline{k^{\d_0\overline{\alpha\b}}_0}
,
\dfrac{\D_0}{g_0}\overline h
\right\rangle
$$
$$
=
\left\langle
k^{\alpha\beta}_\zeta h
,
\D_0
\right\rangle
-
\dfrac{k^{\alpha\beta}_\zeta (0)}{\D_0(0)}
\left\langle
\overline{k^{\d_0\overline{\alpha\b}}_0}
,
\overline h
\right\rangle
$$
by the Direct Cauchy Theorem and by the reproducing kernel property
$$
=
\dfrac
{k^{\alpha\beta}_\zeta(0) h(0)}
{\D_0(0)}
-
\dfrac{k^{\alpha\beta}_\zeta (0)}{\D_0(0)}
h(0)
=0.
$$
Now for $h\in H^2(\overline{\nu_0}\alpha\beta)$
$$
\left\langle
h,
k^{\alpha\beta}_\zeta\overline{g_0}
-
\dfrac{k^{\alpha\beta}_\zeta (0)}{\D_0(0)}
\dfrac{\D_0}{g_0}\overline{k^{\d_0\overline{\alpha\b}}_0}
\right\rangle
$$
$$
=
\left\langle
g_0 h,
k^{\alpha\beta}_\zeta
\right\rangle
-
\dfrac{k^{\alpha\beta}_\zeta (0)}{\D_0(0)}
\left\langle
h g_0 k^{\d_0\overline{\alpha\b}},
\D_0
\right\rangle
$$
$$
=
g_0(\z) h(\z)
-
\dfrac{k^{\alpha\beta}_\zeta (0)}{\D_0(0)}
\dfrac
{h(0) g_0(0) k^{\d_0\overline{\alpha\b}}(0)}
{\D_0(0)}
=
g_0(\z) h(\z)
=\left\langle h, \overline{g_0(\z)}k^{\overline{\nu_0}\alpha\beta}_\zeta\right\rangle .
$$
\eqref{240908_01} follows. Similarly one can verify that the second entry in \eqref{240908-02}
$$
\left(
\overline{w(\zeta)} \dfrac{k^\a_0}{k^\a_0(0)}k^\alpha_\zeta(0)-\overline{w(\zeta)} k^\alpha_\zeta
\right)
\overline{g_0}
-\overline w
\dfrac{k^{\alpha\beta}_\zeta (0)}{\D_0(0)}
\dfrac{\D_0}{g_0}\overline{k^{\d_0\overline{\alpha\b}}_0}
=
-\overline{g_0(\z)}\ \overline{w(\z)}k^{\overline{\nu_0}\alpha}_\zeta.
$$
This completes the proof of the lemma.
\end{proof}
\begin{corollary}\label{250112-01}
Multiplication by $\overline{g_0}$ maps $($isometrically$)$
$$
\sum\limits_k
\left(
\begin{bmatrix}
1 & w \\
\overline{w} & 1
\end{bmatrix}
\begin{bmatrix}
0
\\
k^\a_0
\end{bmatrix}
\dfrac{k^\a_{\z_k}(0)}{k^\a_0(0)}
\ \overline{w(\zeta_k)}
\oplus
K_{\zeta_k}^{\alpha, \beta}
\right)
x_k
$$
to
$$
\sum\limits_k
\left(
K_{\zeta_k}^{\overline{\nu_0}\alpha, \beta}
\ \overline{g_0(\z_k)}
\oplus
\begin{bmatrix}
1 & w \\
\overline{w} & 1
\end{bmatrix}
\begin{bmatrix}
\dfrac{\D_0}{g_0}\overline{k^{\d_0\overline{\alpha\b}}_0}
\\
0
\end{bmatrix}
\dfrac{k^{\a\b}_{\z_k}(0)}{\D_0(0)}
\right)x_k
.
$$
\end{corollary}
%%% This allows to define the family of isometries in terms of the data only.
\begin{corollary}\label{250112-02}
By taking squares of the norms in Corollary \ref{250112-01}, we obtain the family of identities
$$
\sum\limits_{k,j}
\overline{x_j}
w(\z_j)
\dfrac{\overline{k^\a_{\z_j}(0)}k^\a_{\z_k}(0)}{k^\a_0(0)}
\ \overline{w(\zeta_k)}
x_k
+
\sum\limits_{k,j}
\overline{x_j}
\left(
k^{\alpha\beta}_{\z_k}(\zeta_j)-w(\zeta_j)\overline{w(\z_k)} k^\alpha_{\z_k}(\zeta_j)
\right)
x_k
$$
$$
=
\sum\limits_{k,j}
\overline{x_j}
\overline{\left(\dfrac{k^{\a\b}_{\z_j}(0)}{\D_0(0)}\right)}
k^{\d_0\overline{\alpha\b}}_0(0)
\dfrac{k^{\a\b}_{\z_k}(0)}{\D_0(0)}
x_k
$$
$$
+
\sum\limits_{k,j}
\overline{x_j}g_0(\z_j)
\left(
k^{\overline{\nu_0}\alpha\beta}_{\z_k}(\zeta_j)-w(\zeta_j)\overline{w(\z_k)} k^{\overline{\nu_0}\alpha}_{\z_k}(\zeta_j)
\right)
\overline{g_0(\z_k)}
x_k .
$$
\end{corollary}
\begin{corollary}
If a $\b$ automorphic function $w$, analytic on $\bbD$  and bounded in modulus by $1$, solves the Nevanlinna-Pick interpolation problem
\be\label{260227_01}
w(\z_k)=w_k,
\ee
then $w_k$ satisfy the family of identities (isometries)
$$
\sum\limits_{k,j}
\overline{x_j}
w_j
\dfrac{\overline{k^\a_{\z_j}(0)}k^\a_{\z_k}(0)}{k^\a_0(0)}
\ \overline{w_k}
x_k
+
\sum\limits_{k,j}
\overline{x_j}
\left(
k^{\alpha\beta}_{\z_k}(\zeta_j)-w_j\overline{w_k} k^\alpha_{\z_k}(\zeta_j)
\right)
x_k
$$
$$
=
\sum\limits_{k,j}
\overline{x_j}
\overline{\left(\dfrac{k^{\a\b}_{\z_j}(0)}{\D_0(0)}\right)}
k^{\d_0\overline{\alpha\b}}_0(0)
\dfrac{k^{\a\b}_{\z_k}(0)}{\D_0(0)}
x_k
$$
$$
+
\sum\limits_{k,j}
\overline{x_j}g_0(\z_j)
\left(
k^{\overline{\nu_0}\alpha\beta}_{\z_k}(\zeta_j)-w_j\overline{w_k} k^{\overline{\nu_0}\alpha}_{\z_k}(\zeta_j)
\right)
\overline{g_0(\z_k)}
x_k .
$$
and this family of isometries can be isometrically imbedded in the family of unitary colligations \eqref{241225-01}
via Corollary \ref{250112-01}.
\end{corollary}
\begin{remark}
We recall that positivity of all forms
$$
\sum\limits_{k,j}
\overline{x_j}
\left(
k^{\alpha\beta}_{\z_k}(\zeta_j)-w_j\overline{w_k} k^\alpha_{\z_k}(\zeta_j)
\right)
x_k \ge 0,
\quad \forall\a\in\G^*
$$
is necessary and sufficient for existence of a
$\b$ automorphic solution of problem \eqref{260227_01} (see \cite{Kup-Yud-1997}).
\end{remark}


\section{Declarations}
The author has no competing interests to declare that are relevant to the content of this article.
%%%% This manuscript does not report data.
Data sharing is not applicable to this article as no new data were created or analyzed in this study.


\begin{thebibliography}{99}

\bibitem{BR} L. de Branges, J. Rovnyak, \emph{Canonical models in quantum scattering theory.}
In: Perturbation Theory and Its Applications in Quantum Mechanics, Wiley, New York, 1966, 359-391.

%%%%\bibitem{AV}
%%%%D. Alpay, V. Vinnikov, \emph{Indefinite Hardy spaces on finite bordered Riemann surfaces},
%%%%J. Funct. Anal. 172 (2000), no. 1, 221--248.
%%%%
%%%%\bibitem{EVY-2019}
%%%%B. Eichinger, T. Vandenboom, P. Yuditskii, \emph{KDV Hierarchy via Abelian Coverings
%%%%and Operator Identities}, Transactions of the American Mathematical Society, Series B
%%%%6 (2019), 1-44.
%%%%
%%%%\bibitem{BV}
%%%%J. Ball, V. Vinnikov, \emph{Zero-pole interpolation for matrix meromorphic functions on a compact
%%%%Riemann surface and a matrix Fay trisecant identity}, Amer. J. Math. 121 (1999), no. 4, 841--888.

%%%\bibitem{BoKh}
%%%V. Bolotnikov, A. Kheifets,  \emph{A higher order analogue of the Caratheodory-Julia theorem},
%%%Journal of Functional Analysis, 237, no. 1 (2006), 350 - 371

%%%%\bibitem{BS}{A. Borichev, M. Sodin, }\textit{Krein's entire functions and Bernstein
%%%%approximation problem}, Illinois Journal of Mathematics,  45, no. 1 (2001), 167--185.

%%%\bibitem{Car}{C. Carath\' eodory, }\textit{Theory of Functions of a Complex Variable},
%%%Engl. Translation, Chelsea Publishing Company, NY, 1960.

%%%%\bibitem{CarGam}{L. Carleson, T. Gamelin, }\textit{Complex Dynamics}, Springer-Verlag, 1992.
%%%%
%%%%\bibitem{EYu}{A. Eremenko, P. Yuditskii, }  {\em Comb functions.} Recent advances in orthogonal
%%%%polynomials, special functions, and their applications, Contemp. Math.,
%%%%{\bf 578} (2012), 99--118.

%%%\bibitem{Fro}
%%%O. Frostman, \emph{Sur les produits de Blaschke}, Kungl. Fysiografiska S\"alskapets I Lund F\"orhandlingar,
%%%12, no.15 (1943), 169 - 182.

\bibitem{Widom71} H.~Widom,
%%%%{\it $H^p$ sections of vector bundles over Riemann surfaces},
%%%%Ann. of Math., {\bf 94} (1971), 304--324.
{\it The maximum principle for multiple-valued analytic functions},
Acta Math. {126} (1971), 63--82.


\bibitem{Pom} Ch. Pommerenke, \textit{On the Green's function of Fuchsian groups,}
Ann. Acad. Sci. Fenn., 2 (1976), 409-427.


%
\bibitem{Hasu} M. Hasumi,
\emph{Hardy Classes on Infintely Connected Riemann Surfaces},
LNM \textbf{1027}, Springer, New York, Berlin, 1983.

\bibitem{SY}
M. Sodin, P. Yuditskii, \emph{Almost Periodic Jacobi Matrices with Homogeneous Spectrum, Infinite Dimensional Jacobi Inversion, and Hardy Spaces of Character-Automorphic Functions}, Journal of Geometric Analysis
7 (1997) No. 1, 387-435.

\bibitem{Kup-Yud-1997}
S. Kupin, P. Yuditskii, \emph{Analogues of the Nehari and Sarason theorems for character-automorphic
functions and some related questions}, Operator Theory: Advances and Applications 95 (1997), 373-390.

%%%%\bibitem{Yud-1997}
%%%%P. Yuditskii, \emph{The Character-Automorphic Nehari Problem: Non-Uniqueness Criterion and some
%%%%Extremal Solutions}, Journal for Analysis and its Applications
%%%%16 (1997) No. 2, 249-261.
%%%%
%%%%\bibitem{Kh-Yud-2019} A. Kheifets, P. Yuditskii, \emph{Martin Functions of Fuchsian Groups
%%%%and Character Automorphic Subspaces of the Hardy Space in the Upper Half Plane}, To appear, 1-39
%%%%

\bibitem{VY}
A. Volberg, P. Yuditskii,
\emph{Kotani-Last problem and Hardy spaces on surfaces of Widom type},
Invent. Math. 197 (2014), no. 3, 683-740.


%%%%\bibitem{Koo}
%%%%P.~Koosis, \emph{The logarithmic integral. {I}}, Cambridge Studies in Advanced
%%%%  Mathematics, vol.~12, Cambridge University Press, Cambridge, 1998, Corrected
%%%%  reprint of the 1988 original.
%%%%
%%%%
%%%%
%%%%\bibitem{Nik}
%%%%N.~K. Nikolskii, \emph{Treatise on the shift operator}, A Series of
%%%%  Comprehensive Studies in Mathematics, Spriger-Verlag, Berlin Heidelberg New
%%%%  York Tokyo, 1986.
%%%%  %\MR{0827223}
%%%%
%%%%\bibitem{Lev}
%%%%B. Ya. Levin, \textit{Majorants in classes of subharmonic functions, II, The relation between
%%%%majorants and conformal mapping, III, The classification of the closed sets on $\bbR$ and the
%%%%representation of the majorants}, Teor. Funktsii Funktsional. Anal. i Prilozhen. 52 (1989),
%%%%3--33; English translation: J. Soviet Math. 52 (1990), 3351--3372.
%%%%
%%%%\bibitem{LKMV}
%%%%M. S. Liv\v sic, N. Kravitsky, A. S. Markus, V. Vinnikov,
%%%%\textit{Theory of Commuting Nonselfadjoint Operators}, Mathematics and Its Applications,
%%%%332, Kluwer Academic, Dordrecht (1995)
%%%%
%%%%\bibitem{Mar}
%%%%V.~Marchenko, \emph{Sturm-Liouville operators and applications}, Birkh\"auser Verlag, Basel, 1986.

%%\bibitem{PF}
%%B. S. Pavlov, S. I. Fedorov, \textit{Shift group and harmonic analysis on a Riemann surface of genus one},
%%Algebra i Analiz, 1:2 (1989), 132--168; Leningrad Math. J., 1:2 (1990), 447--490


%%%%\bibitem{Kh}
%%%%A. Kheifets, \emph{Abstract interpolation problem and some applications}, 
%%%%Lecture notes given in framework of Holomorphic Spaces semester, fall 1995, 
%%%%in: Holomorphic Spaces (S. Axler, J. McCarthy, D. Sarason editors), MSRI Publications, 
%%%%33 (1998) 351-381, Berkeley, California.
%%%%
%%%%\bibitem{KhCauchy}
%%%%A. Kheifets, \emph{Cauchy Integral Formula for Fuchsian Groups},
%%%%Complex Analysis and Operator Theory, Volume 19, Issue 4 (2025), Article 71
%%%%
%%%%\bibitem{KhCauchy2}
%%%%A. Kheifets, \emph{Cauchy Integral Formula for Fuchsian Groups. II},
%%%%Complex Analysis and Operator Theory, Volume , Issue  (2025), Article 



%%%%\bibitem{VYis}
%%%%A. Volberg and P. Yuditskii, \textit{
%%%%On the inverse scattering problem for Jacobi matrices with the spectrum on an interval,
%%%%a finite system of Intervals or a Cantor set of positive length},
%%%%Communications in Mathematical Physics, {\bf 226} (2002), 567--605.
%%%%
%%%%\bibitem{VY}
%%%%A. Volberg and P. Yuditskii, \textit{Kotani-Last problem and Hardy spaces on
%%%%surfaces of Widom type}. Invent. Math., 197 (2014), No. 3, 683-740.
%%%%
%%%%\bibitem{VYM}
%%%%A. Volberg and P. Yuditskii, \textit{Mean type of functions of bounded characteristic
%%%%and Martin functions in Denjoy domains}, Adv. in Math.,  290 (2016), 860--887.

%
%%%\bibitem{Widom} H.~Widom,
%%%{\it Extremal polynomials associated with a system of curves in the complex plane},
%%%Adv. in Math., {\bf 3} (1969), 127--232.





%%%%\bibitem{You}
%%%%J. You,
%%%%\textit{Quantitative almost reducibility and its applications},
%%%%Proc. Int. Cong. of Math. - 2018, Rio de Janeiro, Vol. 2, 2107--2128.


\end{thebibliography}

\bigskip

\bigskip



A. Kheifets, Department of Mathematics and Statistics, University of Massachusetts Lowell, One University Ave.,
Lowell, MA 01854,USA

\emph{E-mail address:} {Alexander\underline{ }Kheifets@uml.edu}




 \end{document} 
